\documentclass[11pt,a4paper]{article}
\usepackage{amsmath, amssymb}
\usepackage[margin=1in]{geometry}
\usepackage{parskip}

\title{Mathematical Modeling of a DC Motor}
\author{}
\date{}

\begin{document}
\maketitle

%--------------------------------------------------------------
\section{Equation of Motion}
%--------------------------------------------------------------

The first-order rotational dynamics of a DC motor with viscous friction are
\begin{equation}\label{eq:ode}
  \frac{d\omega}{dt} = \frac{T(t)}{I} - \frac{C}{I}\,\omega,
\end{equation}
where $\omega(t)$ is the angular velocity, $T(t)$ is the applied torque,
$I$ is the moment of inertia, and $C$ is the viscous damping coefficient.

%--------------------------------------------------------------
\section{Torque Input (PWM Pulse)}
%--------------------------------------------------------------

The torque is modelled as a rectangular pulse of duration $t_n$ derived from
a PWM signal:
\begin{equation}\label{eq:torque}
  T(t) = \alpha\, P\bigl[u(t) - u(t - t_n)\bigr],
\end{equation}
where $\alpha$ is a torque constant (relating PWM to torque),
$P$ is the PWM duty-cycle amplitude, and $u(t)$ is the Heaviside unit-step
function.

%--------------------------------------------------------------
\section{Laplace-Domain Solution for $\omega$}
%--------------------------------------------------------------

Taking the Laplace transform of \eqref{eq:ode} with initial condition
$\omega(0)=\omega_0$:
\[
  s\,\Omega(s) - \omega_0 + \frac{C}{I}\,\Omega(s)
  = \frac{\alpha P}{I}\;\frac{1 - e^{-t_n s}}{s}.
\]
Solving for $\Omega(s)$:
\begin{equation}\label{eq:Omega_s}
  \Omega(s)
  = \frac{\omega_0}{s + \frac{C}{I}}
  + \frac{\alpha P}{I}\;\frac{1 - e^{-t_n s}}{s\!\left(s + \frac{C}{I}\right)}.
\end{equation}

Using partial fractions,
\[
  \frac{1}{s\!\left(s+\frac{C}{I}\right)}
  = \frac{I}{C}\left(\frac{1}{s} - \frac{1}{s+\frac{C}{I}}\right).
\]

%--------------------------------------------------------------
\section{Time-Domain Solution for $\omega(t)$}
%--------------------------------------------------------------

Inverting \eqref{eq:Omega_s} gives
\begin{equation}\label{eq:omega_t}
\boxed{
  \omega(t)
  = \omega_0\, e^{-\frac{C}{I}t}
  + \frac{\alpha P}{C}
    \Bigl[\bigl(1 - e^{-\frac{C}{I}t}\bigr)
    - u(t-t_n)\bigl(1 - e^{-\frac{C}{I}(t-t_n)}\bigr)\Bigr].
}
\end{equation}

\textbf{Equivalently} (factoring as in the handwritten notes with $\alpha P/I$
out front):
\[
  \omega(t)
  = \omega_0\, e^{-\frac{C}{I}t}
  + \frac{\alpha P}{I}\left[
      \frac{I}{C} - \frac{I}{C}\,e^{-\frac{C}{I}t}
      - \frac{I}{C}\,u(t-t_n)
      + \frac{I}{C}\,e^{-\frac{C}{I}(t-t_n)}\,u(t-t_n)
    \right].
\]

%--------------------------------------------------------------
\section{Angular Position $\theta(t)$ by Integration}
%--------------------------------------------------------------

Integrating $\omega(t)$ with $\theta(0)=\theta_0$:
\[
  \theta(t) = \theta_0 + \int_0^{t} \omega(\tau)\,d\tau.
\]

Carrying out the integration term by term:

\begin{equation}\label{eq:theta_t}
\boxed{
\begin{aligned}
  \theta(t)
  &= \theta_0
   + \omega_0\,\frac{I}{C}\!\left(1 - e^{-\frac{C}{I}t}\right) \\[6pt]
  &\quad + \frac{\alpha P}{C}\left[
      t + \frac{I}{C}\,e^{-\frac{C}{I}t} - \frac{I}{C}
    \right] \\[6pt]
  &\quad - \frac{\alpha P}{C}\;u(t-t_n)\left[
      (t - t_n) + \frac{I}{C}\,e^{-\frac{C}{I}(t-t_n)} - \frac{I}{C}
    \right].
\end{aligned}
}
\end{equation}

Or equivalently, factoring $\alpha P / I$ out front:
\[
\begin{aligned}
  \theta(t)
  &= \theta_0
   + \omega_0\,\frac{I}{C}\!\left(1 - e^{-\frac{C}{I}t}\right) \\[6pt]
  &\quad + \frac{\alpha P}{I}\left[
      \frac{I}{C}\,t
      + \frac{I^2}{C^2}\,e^{-\frac{C}{I}t}
      - \frac{I^2}{C^2}
    \right] \\[6pt]
  &\quad - \frac{\alpha P}{I}\;u(t-t_n)\left[
      \frac{I}{C}(t-t_n)
      + \frac{I^2}{C^2}\,e^{-\frac{C}{I}(t-t_n)}
      - \frac{I^2}{C^2}
    \right].
\end{aligned}
\]

%--------------------------------------------------------------
\section{Hall Sensor Observation Model}
%--------------------------------------------------------------

The rotor position $\theta(t)$ is observed indirectly through a binary Hall
sensor.  For a 4-pole motor the sensor outputs 1 when the rotor angle lies
inside a window of half-width $k$ centred on each multiple of $\pi/2$:
\begin{equation}\label{eq:hall_hard}
  y(t) = \mathbf{1}\!\left[
    \exists\, n\in\mathbb{Z}:\;
    \bigl|\theta(t) - n\tfrac{\pi}{2}\bigr| < k
  \right].
\end{equation}
Equivalently, defining the distance from the nearest window centre
\begin{equation}
  d(t) = \min\!\bigl(\theta(t)\bmod\tfrac{\pi}{2},\;
                     \tfrac{\pi}{2} - \theta(t)\bmod\tfrac{\pi}{2}\bigr),
\end{equation}
the hard indicator is $y(t)=\mathbf{1}[d(t)<k]$.

\subsection{Hard indicator — pole-tracking implementation}

Because $\theta(t)$ is monotonically increasing during motor operation, the
implementation tracks a \emph{pole counter} $p$ that advances whenever
$\theta$ exits the current window from above, rather than recomputing the
modulus at every step.  For each discrete sample $\theta_i$:

\begin{enumerate}
  \item \textbf{Advance pole}: while $\theta_i \ge p\,\dfrac{\pi}{2} + k$,
        set $p \leftarrow p + 1$.
        (Handles the case where a single time step jumps over one or more
        complete windows at high angular speed.)
  \item \textbf{Classify}:
        \[
          y_i =
          \begin{cases}
            1 & \text{if } \theta_i \ge p\,\dfrac{\pi}{2} - k, \\
            0 & \text{otherwise.}
          \end{cases}
        \]
\end{enumerate}

This runs in $O(N)$ total (amortised $O(1)$ per sample) and is implemented
in \texttt{omega2signal($\theta_i$, $p$, $k$)} which returns both $y_i$ and
the updated pole $p$ to carry forward.

\textbf{Transition detection.}
The binary signal changes from 0 to 1 (rising edge) or 1 to 0 (falling
edge) whenever $y_i \neq y_{i-1}$.  A rising edge corresponds to $\theta$
entering a window from below; a falling edge to $\theta$ leaving a window
from above (or jumping directly into the next window at high speed, in which
case consecutive rising/falling edges may be skipped and only the net
transition is recorded).

\subsection{Soft (differentiable) approximation}
To enable gradient-based optimisation, the hard indicator is replaced by a
smooth sigmoid:
\begin{equation}\label{eq:hall_soft}
  \hat{p}(t;\,\mathbf{x},s)
  = \sigma\!\bigl(s\,(k - d(t))\bigr), \qquad
  \sigma(z) = \frac{1}{1+e^{-z}},
\end{equation}
where $s>0$ is a \emph{sharpness} hyperparameter.
As $s\to\infty$, $\hat{p}\to y$ pointwise.

%--------------------------------------------------------------
\section{Parameter Identification Problem}
%--------------------------------------------------------------

\subsection{Parameters}
The five scalar parameters to be identified are
\[
  \mathbf{x} = (a_I,\; c_I,\; \omega_0,\; \theta_0,\; k),
\]
where $a_I = \alpha/I$, $c_I = C/I$, $\omega_0$ is the initial angular
velocity, $\theta_0$ is the initial angle, and $k$ is the Hall window
half-width.

\subsection{Data alignment}
The raw Arduino log (timestamps in microseconds, binary Hall state) is
filtered to the interval $[2.26\,\text{s},\,14\,\text{s}]$ and re-zeroed
so that $t=0$ corresponds to PWM onset.  The PWM pulse runs from $t=0$ to
$t=t_n=10\,\text{s}$.

%--------------------------------------------------------------
\section{Loss Functions Evaluated}
%--------------------------------------------------------------

Let $\{(t_i, y_i)\}_{i=1}^{N}$ be the filtered sensor samples.

\subsection{Binary MSE (\texttt{score})}
\begin{equation}
  \mathcal{L}_{\text{score}}(\mathbf{x})
  = \frac{1}{N}\sum_{i=1}^{N}\!\bigl(\hat{y}_i - y_i\bigr)^2,
  \qquad \hat{y}_i \in \{0,1\}.
\end{equation}
Non-differentiable; many degenerate local minima at $\approx 0.36$.

\subsection{Transition-timestamp MSE (\texttt{theta\_function})}
Let $\{t^{\text{pred}}_j\}$ and $\{t^{\text{act}}_j\}$ be the predicted and
actual times of 0/1 transitions.  With $n=\min(|\text{pred}|,|\text{act}|)$:
\begin{equation}
  \mathcal{L}_{\text{time}}(\mathbf{x})
  = \frac{1}{n}\sum_{j=1}^{n}\!\bigl(t^{\text{pred}}_j - t^{\text{act}}_j\bigr)^2
  + \left(\frac{|\text{pred}|-|\text{act}|}{|\text{act}|}\right)^{\!2}.
\end{equation}
Smoother than binary MSE but still has a flat valley where total transition
count matches but timing is globally shifted; best observed score $\approx 0.46$.

\subsection{Transition-interval MSE (\texttt{interval\_loss})}
\begin{equation}
  \mathcal{L}_{\text{interval}}(\mathbf{x})
  = \frac{1}{n-1}\sum_{j=1}^{n-1}
    \!\bigl(\Delta t^{\text{pred}}_j - \Delta t^{\text{act}}_j\bigr)^2
  + \left(\frac{|\text{pred}|-|\text{act}|}{|\text{act}|}\right)^{\!2},
\end{equation}
where $\Delta t_j = t_{j+1}-t_j$ is the inter-transition interval (a
direct proxy for instantaneous angular velocity $\omega \approx \pi/(2\Delta
t)$).  Immune to global phase offset; best observed score $\approx 0.46$.

\subsection{Binary cross-entropy with soft Hall window (\texttt{bce\_loss})}
Using the soft approximation \eqref{eq:hall_soft}:
\begin{equation}\label{eq:bce}
  \mathcal{L}_{\text{BCE}}(\mathbf{x};\,s)
  = -\frac{1}{N}\sum_{i=1}^{N}\!
    \Bigl[y_i\log\hat{p}_i + (1-y_i)\log(1-\hat{p}_i)\Bigr].
\end{equation}
Fully differentiable; correct Bernoulli likelihood for the observation model.
Dense gradient signal (every sample contributes, not just transitions).

%--------------------------------------------------------------
\section{Optimisation Methods Evaluated}
%--------------------------------------------------------------

\subsection{L-BFGS-B (gradient-based)}
Applied to $\mathcal{L}_{\text{score}}$ initially; failed because the loss is
non-differentiable.  After switching to $\mathcal{L}_{\text{BCE}}$, 200
random restarts with L-BFGS-B achieved score $\approx 0.45$.  Fast per
restart but gets trapped in local minima.

\subsection{CMA-ES (\texttt{cmaes} library)}
Gradient-free covariance-matrix adaptation.  Population 1000, 200
generations, parameters normalised to $[0,1]$.  Applied to
$\mathcal{L}_{\text{score}}$: best score $\approx 0.363$ (degenerate
high-speed / small-$k$ solution); applied to $\mathcal{L}_{\text{BCE}}$
(sharpness 50): score $\approx 0.59$.

\subsection{Differential Evolution (\texttt{scipy.optimize.differential\_evolution})}
Stochastic population-based global search with optional L-BFGS-B polish.
Applied to $\mathcal{L}_{\text{BCE}}$:
\begin{itemize}
  \item \texttt{popsize=15}: BCE $=1.52$, score $=0.214$.
  \item \texttt{popsize=100}: BCE $=0.664$, score $=0.405$.
\end{itemize}
Revealed an \emph{inversion}: lower BCE does not imply lower binary MSE,
because the optimizer exploits the soft window to spread uncertainty
($\hat{p}\approx 0.5$) rather than making correct hard predictions.

\subsection{Sharpness annealing (forward: $s=2\to50$)}
Sequential DE stages with increasing sharpness:
\begin{center}
\begin{tabular}{ccc}
\hline
$s$ & BCE & Score \\ \hline
2  & 0.554 & 0.244 \\
5  & 0.438 & 0.189 \\
15 & 0.633 & 0.214 \\
50 & 1.506 & 0.216 \\ \hline
\end{tabular}
\end{center}
Best score achieved: $\mathbf{0.189}$ at $s=5$.

\subsection{Sharpness annealing (reverse: $s=5\to1$)}
Schedule $s\in\{5,4,3,2,1\}$, each stage warm-started from the previous
best (\texttt{popsize=15}, \texttt{maxiter=500}):
\begin{center}
\begin{tabular}{ccc}
\hline
$s$ & BCE & Score \\ \hline
5 & 0.4918 & 0.2189 \\
4 & 0.5033 & \textbf{0.1715} \\
3 & 0.5295 & 0.1977 \\
2 & 0.5440 & 0.2428 \\
1 & 0.6046 & 0.3777 \\ \hline
\end{tabular}
\end{center}
Best score achieved: $\mathbf{0.1715}$ at $s=4$.

\subsection{Optimal sharpness analysis}

The observed optimal $s^*\approx 4$ can be understood from the sigmoid
transition width.  The 10\%--90\% width of $\sigma(s\,(k-d))$ around the
window edge $d=k$ is $\Delta d_{90} = \frac{\ln 9}{s} \approx \frac{2.2}{s}$.
At $s=4$, $\Delta d_{90}\approx 0.55\,\text{rad}$, which is comparable to the
window half-width $k\approx 0.341\,\text{rad}$.

\begin{itemize}
  \item $s \gg 1/k$: nearly hard threshold --- optimizer exploits the sharp
        boundary by placing $\theta$ exactly at the edge, giving high
        confidence on the wrong side for a fraction of samples; BCE penalises
        these heavily, so BCE increases but score worsens.
  \item $s \ll 1/k$: very wide soft window --- the gradient is weak
        everywhere; the optimizer settles for $\hat{p}\approx 0.5$
        (uniform uncertainty), BCE approaches $\ln 2\approx 0.693$, and
        score is near-random.
  \item $s \approx 1/k$ (here $s\approx 3$--$5$): the sigmoid width matches
        the physical window width, giving the tightest alignment between the
        soft loss landscape and the binary MSE we ultimately care about.
\end{itemize}

\subsection{Adaptive sharpness: $s = 1/k$}

Since $k$ is an unknown being optimised, fixing $s$ manually requires
re-tuning whenever $k$ changes.  The natural solution is to couple
\[
  s(\mathbf{x}) = \frac{1}{k},
\]
computed from the current parameter estimate at every function evaluation.
This keeps the sigmoid width matched to the window width throughout
optimisation with no additional hyperparameter.

\textbf{Result (single-stage DE, \texttt{popsize=15}, \texttt{maxiter=1000}):}
BCE $= 0.459$, score $= 0.230$.

Compared with the manual sharpness annealing best of $0.1715$, the single-stage
adaptive run is slightly worse.  The annealing schedule provided implicit
warm-starting that is lost here; combining adaptive sharpness with
warm-starting across stages is a natural next step.

%--------------------------------------------------------------
\subsection{3-sensor BCE optimisation (\texttt{bce\_loss\_3sensor})}
%--------------------------------------------------------------

A cosine-based soft model is used for the 3-sensor case, eliminating $k$
entirely.  With electrical angle $\varphi = 2\theta$:
\begin{align}
  \hat{p}_A &= \sigma\!\bigl(s\cos\varphi\bigr), \notag\\
  \hat{p}_B &= \sigma\!\bigl(s\cos(\varphi - \tfrac{2\pi}{3})\bigr), \notag\\
  \hat{p}_C &= \sigma\!\bigl(s\cos(\varphi - \tfrac{4\pi}{3})\bigr). \notag
\end{align}
$\cos = +1$ at the window centre, $-1$ at the gap centre, and $0$ at the
window edge (50\% duty cycle), so the sharpness $s$ acts analogously to
$s/k$ in the single-sensor model.  The optimised parameters reduce to
$\mathbf{x} = (a_I, c_I, \omega_0, \theta_0)$ (no $k$), and the loss is:
\begin{equation}
  \mathcal{L}_{\text{3s}}(\mathbf{x};\,s)
  = -\frac{1}{N}\sum_{i=1}^{N}
    \bigl[A_i\log\hat{p}_{A,i} + (1-A_i)\log(1-\hat{p}_{A,i})
         + B_i(\ldots) + C_i(\ldots)\bigr].
\end{equation}

\textbf{Diagnosis of failure.}  At steady state ($\omega_{ss}\approx 586$
rad/s) with $\Delta t \approx 0.96\,\text{ms}$, the 3-bit Hall state changes
every $\approx 1.07$ samples.  Any error in the implied $\omega_{ss}$ causes
rapid phase accumulation: a 1\% error places the predicted state one step
wrong after $\approx 100$ samples, giving BCE $\gg 2.08$ (the 3-sensor random
baseline).  \textbf{Result: BCE $=4.16$, score $=0.439$.}

A \emph{constrained} 3-sensor variant (\texttt{bce\_loss\_3sensor\_constrained})
fixes $\omega_{ss}$ from the 3-sensor measurement and optimises only
$(c_I,\omega_0,\theta_0)$, with $a_I = \omega_{ss}\,c_I/P$.
\textbf{Result: BCE $=4.24$, score $=0.444$} --- marginally worse, confirming
that the 4-parameter free space was not the bottleneck; the binary cosine model
itself is too phase-sensitive at this speed.

%--------------------------------------------------------------
\subsection{Hybrid approach: 3-sensor $\omega_{ss}$ + single-sensor BCE}
%--------------------------------------------------------------

Since the 3-sensor measurement of $\omega_{ss}$ is accurate (no pole-skipping,
12 transitions/revolution), while the 3-sensor BCE optimisation fails due to
extreme phase sensitivity, the natural hybrid is:

\begin{enumerate}
  \item Compute $\omega_{ss}$ from 3-sensor state transitions via
        $\omega_{ss} = \pi / (6\,\overline{\Delta t}_{\text{state}})$.
  \item Estimate $k$ from the single-sensor duty cycle in the steady-state
        window: $k = \overline{y}_{ss}\,\pi/4$ (unaffected by pole-skipping).
  \item Optimise the full 5-parameter single-sensor BCE with all parameters
        free; use the 3-sensor $\omega_{ss}$ only to seed the first population
        row via $a_{I,\text{seed}} = \omega_{ss}\,c_{I,\text{seed}}/P$.
\end{enumerate}

\textbf{Result:} BCE $=0.727$, score $=0.476$, implied $\omega_{ss}=1047$ rad/s.
The optimizer found a spurious solution with $\omega_0=771$ and $c_I=6$ ---
the motor reaching steady state almost instantly, with the incorrect
$a_I/c_I$ ratio.

%--------------------------------------------------------------
\subsection{Sharpness as a free parameter}
%--------------------------------------------------------------

A wrapper \texttt{\_bce\_with\_sharpness} expands the parameter vector to
$\mathbf{x} = (a_I,c_I,\omega_0,\theta_0,k,s)$ and optimises sharpness
jointly.  An informed seed is placed at $s_{\text{seed}} = 1/k_{\text{meas}}$.

\textbf{Result:} score $=0.451$, $s^*=6.2$, implied $\omega_{ss}=1047$ rad/s
--- sharpness landed near the expected range but the speed degeneracy
remained.

%--------------------------------------------------------------
\subsection{Transient-only optimisation}
%--------------------------------------------------------------

\textbf{Motivation.}  The steady-state plateau ($t \in [3,10]\,\text{s}$,
$\approx 60\%$ of all samples) constrains only the ratio $a_I/c_I = \omega_{ss}/P$.
The ramp-up ($t \in [0,3]\,\text{s}$) and coast-down ($t \in [10,14]\,\text{s}$)
both exhibit curvature that independently constrains $c_I = 1/\tau$, so
removing the plateau forces the optimizer to fit the transient shape.

\textbf{Implementation.}  The data array is masked to keep only
$t \leq 3\,\text{s}$ and $t \geq t_n$.  The split index $p_1$ is
recomputed as \texttt{searchsorted}$(t_{\text{sub}}, t_n)$ so that
\texttt{\_theta} still applies the correct analytical formula on each segment.

\textbf{Result:} BCE $=0.678$, score $=0.451$.  The recovered
$\omega_{ss,\text{implied}} = 9492$ rad/s and $k=0.13$, indicating a
degenerate solution with tiny window and near-constant $\hat{p}\approx 0.5$.

\textbf{Lesson.}  The steady-state plateau is \emph{necessary}: it anchors
the Hall transition frequency (hence $a_I/c_I$).  Without it, the BCE
loss can be nearly minimised by predicting $\hat{p}\approx 0.5$ everywhere
(BCE $\approx\ln 2 = 0.693$), making the transient shape irrelevant.

%--------------------------------------------------------------
\subsection{Omega\_ss soft penalty (\texttt{\_bce\_omega\_penalty})}
%--------------------------------------------------------------

\textbf{Motivation.}  The hard constraint $a_I = \omega_{ss}\,c_I/P$ couples
$a_I$ and $c_I$ and prevents the transient from independently constraining
$c_I$.  A soft penalty preserves all data and all degrees of freedom:
\begin{equation}\label{eq:penalty}
  \mathcal{L}_{\text{total}}(\mathbf{x};\,s,\lambda)
  = \mathcal{L}_{\text{BCE}}(\mathbf{x};\,s)
  + \lambda\!\left(\frac{a_I P}{c_I} - \omega_{ss}\right)^{\!2}.
\end{equation}

The parameter vector is reduced back to
$\mathbf{x}=(a_I,c_I,\omega_0,\theta_0,k)$ with $s$ fixed per stage (not
free), and $\omega_{ss}$ supplied from the 3-sensor measurement.

\subsubsection*{Single-stage with free sharpness ($\lambda=10^{-4}$)}

With sharpness also free (6 parameters), the optimizer found
$s^*=0.62$ and $k=0.002$, i.e.\ $\hat{p}\approx 0.5$ everywhere.
BCE $=0.685$, score $=0.434$; $\omega_{ss}$ matched the measurement but
$a_I$, $c_I$ were wrong.  Conclusion: the free sharpness allows the optimizer
to trivially minimise BCE by making the window nearly invisible.

\subsubsection*{Sharpness annealing ($\lambda=10^{-4}$, schedule $s\in\{2,3,4,5\}$)}

With sharpness fixed per stage (5-parameter optimisation) and warm-starts:
\begin{center}
\begin{tabular}{cccc}
\hline
$s$ & Total loss & Score & Implied $\omega_{ss}$ (rad/s) \\ \hline
2 & 0.699 & 0.468 & 543 \\
3 & 0.721 & 0.466 & 543 \\
4 & 0.753 & 0.467 & 543 \\
5 & 0.754 & \textbf{0.244} & 586 \\ \hline
\end{tabular}
\end{center}
At $s=5$ the sharpness term dominates and the penalty is overcome, allowing
the implied $\omega_{ss}$ to drift from 543 to 586 rad/s.  Score drops to
0.244, which is explained by steady-state dominance: matching the Hall
frequency over 60\% of the data (steady state) alone yields
score $\approx 0.6 \times 0 + 0.4 \times 0.5 = 0.20$.
Recovered $\omega_0=947$ rad/s confirms the optimizer found a fast-ramp
degenerate solution.

\subsubsection*{Tightened $\omega_0$ bound and stronger penalty ($\lambda=10^{-3}$)}

Two additional constraints are physically motivated:
\begin{itemize}
  \item $\omega_0 \in [0,10]\,\text{rad/s}$: the motor is always at rest
        before PWM is applied, eliminating fast-ramp degenerate solutions.
  \item $\lambda = 10^{-3}$: at 1\% $\omega_{ss}$ error the penalty
        contributes $10^{-3}\times(5.4)^2\approx 0.029$, non-negligible
        relative to a BCE of $\approx 0.75$.
\end{itemize}

\textbf{Result:} score $=0.468$, implied $\omega_{ss}=540.6$ rad/s (within
0.03\% of measurement), $\omega_0=1.5$ rad/s.  However, $a_I=0.280$,
$c_I=0.803$ --- both an order of magnitude below the true values.  The
$a_I/c_I$ ratio is correct (hence $\omega_{ss}$ matches), but their absolute
scale is not.

\textbf{Root cause: $a_I$--$c_I$ scale degeneracy.}  For a given
$\omega_{ss}=a_IP/c_I$, any pair $(\lambda a_I, \lambda c_I)$ for
$\lambda>0$ produces the same asymptotic speed.  Distinguishing them
requires the transient shape; but the binary BCE loss over a long steady-state
window provides very weak gradient signal for the transient parameters once
$\omega_{ss}$ is pinned.  This degeneracy is a fundamental limitation of
the single Hall sensor approach under this loss formulation.

%--------------------------------------------------------------
\section{Analytical Parameter Estimation Pipeline}
%--------------------------------------------------------------

\subsection{Data alignment bug and fix}

A critical bug was discovered in \texttt{data.py}: \texttt{load\_data()} called
\texttt{t = t - t[0]} \emph{before} filtering to $t\geq 2.26\,\text{s}$.
After filtering, \texttt{t[0]}\,$\approx 2.26$ rather than $0$, so every
subsequent computation treated the time axis as starting at the pre-PWM
epoch rather than at PWM onset.

\textbf{Effect on coast-down estimation.}  The function
\texttt{estimate\_c\_I\_coastdown} selects transitions with $t > t_n = 10\,\text{s}$.
With the bug, $t$ ran from $2.26$ to $14\,\text{s}$, so the ``coast-down''
window $t > 10$ captured $\approx 1.74\,\text{s}$ of actual deceleration
preceded by $\approx 1.74\,\text{s}$ of steady-state data
(physical $t \in [10, 12.26]\,\text{s}$).
The log-linear slope was $\hat{c}_I \approx 0.358\,\text{s}^{-1}$ (true value
$2.0\,\text{s}^{-1}$), an error of $5.6\times$.

\textbf{Effect on optimisation.}  All previous DE runs were fitting a
model with $t$ misaligned by $2.26\,\text{s}$.  The optimizer compensated with
$\omega_0 \sim 944\,\text{rad/s}$ and large $c_I$, giving a spurious
``fast-ramp'' solution.

\textbf{Fix.}  A second re-zeroing step is added \emph{after} filtering in
both \texttt{load\_data()} and \texttt{load\_data\_3sensor()}:
\begin{verbatim}
index = t >= 2.26
t = t[index];  w = w[index]
t = t - t[0]   # re-zero to PWM onset
\end{verbatim}
After the fix: $t=0$ corresponds to PWM onset; the coast-down correctly
begins at $t = t_n = 10\,\text{s}$.

\subsection{Coast-down $c_I$ estimation}

During coast-down ($t > t_n$, no applied torque):
\[
  \omega(t) = \omega_{ss}\,e^{-c_I(t-t_n)},
\]
so consecutive Hall transition intervals $\Delta\tau_j$ satisfy
\[
  \Delta\tau_j \approx \frac{\pi/4}{\omega(t_{\text{mid},j})}
  = \frac{\pi}{4\,\omega_{ss}}\,e^{+c_I(t_{\text{mid},j}-t_n)}.
\]
Taking logarithms:
\[
  \ln\Delta\tau_j = \underbrace{\ln\!\frac{\pi}{4\,\omega_{ss}}}_{\text{const}}
                   + c_I\,(t_{\text{mid},j} - t_n).
\]
A linear regression of $\ln\Delta\tau$ vs.\ $(t_{\text{mid}}-t_n)$ yields
the slope $\hat{c}_I$ directly.

\textbf{Result (after data-alignment fix):} $\hat{c}_I = 2.024\,\text{s}^{-1}$
vs.\ true $c_I = 2.0\,\text{s}^{-1}$ (1.2\% error).

\subsection{3-sensor quantization floor at high speed}

The 3-sensor steady-state estimator computes
$\omega_{ss} = \pi / (6\,\overline{\Delta t}_{\text{state}})$,
where $\overline{\Delta t}_{\text{state}}$ is the mean interval between
consecutive state changes.  At the operating speed:
\[
  \overline{\Delta t}_{\text{state}}
  = \frac{\pi/6}{\omega_{ss}}
  = \frac{\pi/6}{586\,\text{rad/s}}
  \approx 0.893\,\text{ms}
  < \Delta t = 0.962\,\text{ms}.
\]
Because the motor changes Hall state \emph{faster than the sampling rate},
every measured interval is quantized to exactly $\Delta t$.  The estimator
therefore returns
\[
  \hat{\omega}_{ss}^{\text{(3-sensor)}}
  = \frac{\pi}{6\,\Delta t}
  = \frac{\pi}{6 \times 0.962\,\text{ms}}
  = 544.3\,\text{rad/s},
\]
regardless of the actual speed.  This is a hard quantization floor, not a
statistical error: the 3-sensor estimator \emph{cannot} produce a value
above $\pi/(6\Delta t)$.

\textbf{Single-sensor comparison.}  The single-sensor transitions correspond
to $\pi/4$ rad, giving intervals of $\approx \pi/(4\times 586) = 1.34\,\text{ms} > \Delta t$.
Each interval is measurable, yielding:
\[
  \hat{\omega}_{ss}^{\text{(single)}} = 586.1\,\text{rad/s}
  \quad\text{(true: }586.2\,\text{rad/s, error }0.02\%).
\]
At this operating point, the single-sensor estimate is more accurate than
the 3-sensor estimate.

\textbf{The quantization floor is a general condition.}  The 3-sensor
estimator fails whenever
\[
  \omega_{ss} > \frac{\pi}{6\,\Delta t} \approx 544\,\text{rad/s}
  \quad (5{,}200\,\text{RPM at this baud rate}).
\]

\subsection{Analytical estimation pipeline}

Given the identified issues, parameters are recovered analytically before
any numerical optimisation:

\begin{enumerate}
  \item \textbf{$\omega_{ss}$}: from single-sensor transitions in the
        steady-state window $[3, 9]\,\text{s}$ via
        $\hat\omega_{ss} = \pi / (4\,\overline{\Delta t})$.
  \item \textbf{$k$}: from the steady-state duty cycle,
        $\hat{k} = \overline{y}_{ss}\,\pi/4$.
  \item \textbf{$c_I$}: from the coast-down log-linear fit (Section above).
  \item \textbf{$a_I$}: derived as $\hat{a}_I = \hat\omega_{ss}\,\hat{c}_I / P$.
  \item \textbf{$\omega_0$}: fixed to $0$ (motor starts from rest).
  \item \textbf{$\theta_0$}: estimated via a three-phase grid search.
\end{enumerate}

\subsubsection*{Phase 1: 1D grid over $\theta_0$}
Since $\theta_0$ is periodic mod $\pi/2$, a grid of 200 points over
$[0,\pi/2)$ evaluates the BCE loss for each $\theta_0$ with all other
parameters fixed at their analytic estimates.
\textbf{Best score: 0.302 at $\theta_0 = 0.118$.}

\subsubsection*{Phase 2: 2D Nelder-Mead over $(a_I, \theta_0)$}
A Nelder-Mead local search refines $a_I$ and $\theta_0$ jointly (all
other parameters fixed).
\textbf{Result: $a_I=0.764$, $\theta_0=-0.331$, score $=0.194$.}

\subsubsection*{Phase 3: Parallel exhaustive 2D grid}
Phase drift sensitivity motivates a dense search: a $1000 \times 200$
grid over $(\omega_{ss} \pm 4\,\text{rad/s}) \times [0, \pi/2)$.
For each $\omega_{ss}$ point, $a_I = \omega_{ss}\,c_I/P$ is derived.
The grid is evaluated in parallel via \texttt{multiprocessing.Pool}.

\textbf{Best result:}
\begin{center}
\begin{tabular}{lcc}
\hline
Parameter & Recovered & True \\ \hline
$a_I$ & 0.764 & 0.755 \\
$c_I$ & 2.024 & 2.000 \\
$\omega_0$ & 0.0 (fixed) & 0.000 \\
$\theta_0$ & 1.084 & 1.000 \\
$k$ & 0.341 & 0.341 \\
$\omega_{ss}$ & 586.23 & 586.24 \\
\hline
Score & \multicolumn{2}{c}{\textbf{0.180}} \\
\hline
\end{tabular}
\end{center}

\subsection{Phase drift from $\omega_{ss}$ error}

At $\omega_{ss} \approx 586\,\text{rad/s}$, even a small relative error
$\epsilon$ in $\omega_{ss}$ accumulates rapidly:
\[
  \Delta\theta(t) \approx \epsilon\,\omega_{ss}\,t.
\]
After $t = 10\,\text{s}$ of steady-state rotation, a 0.017\% error
($\approx 0.1\,\text{rad/s}$) gives $\Delta\theta \approx 1\,\text{rad}$,
placing approximately 50\% of predictions in the wrong binary state.
This is why the grid search over $\omega_{ss}$ (Phase 3) is necessary
even when $\hat\omega_{ss}$ is already accurate to 0.02\%.

\subsection{Lower-PWM recommendation for reliable identification}

All quantization and phase-accumulation problems become tractable at lower
motor speed.  The 3-sensor quantization floor is $\pi/(6\Delta t) \approx 544\,\text{rad/s}$.
To guarantee that 3-sensor intervals are measurable:
\[
  \omega_{ss} < \frac{\pi}{6\,\Delta t} = 544\,\text{rad/s}
  \;\Longleftrightarrow\;
  P_{\text{scale}} \lesssim 0.93.
\]
For comfortable margin (3-sensor resolution $\geq 2\Delta t$):
\[
  \omega_{ss} < \frac{\pi}{12\,\Delta t} = 272\,\text{rad/s}
  \;\Longleftrightarrow\;
  P_{\text{scale}} \approx 0.46.
\]
At $\omega_{ss} = 272\,\text{rad/s}$, state-change intervals are $\approx 3.86\,\text{ms}
= 4\,\Delta t$: well-resolved by both sensors.  Phase drift over 10\,s at
a 0.1\,rad/s error is only $\approx 1\,\text{rad}$, the same absolute amount
but occurring over a smaller fraction of the rotation cycle.
Additionally, the coast-down takes longer and provides more transitions for
the log-linear $c_I$ fit.

The \texttt{generate\_data.py} script now accepts a \texttt{P\_SCALE}
parameter (default $1.0$, set to $\approx 0.46$ for the low-speed
calibration run) to regenerate synthetic data at any desired operating point.

%--------------------------------------------------------------
\section{Computational Implementation}
%--------------------------------------------------------------

\subsection{Vectorised \texttt{score} and \texttt{pole=0} fix}

The original \texttt{score} used a Python \texttt{while} loop for
pole-tracking over $N\approx11{,}000$ samples, and initialised
\texttt{pole=1}, missing the Hall window at $\theta\in[-k,k]$ (the $n=0$
pole).  Both issues are fixed by replacing the loop with three vectorised
NumPy operations, exploiting the equivalence
\[
  y(t) = \mathbf{1}[d(t) < k], \qquad
  d(t) = \min\!\bigl(\theta\bmod\tfrac{\pi}{2},\;
                     \tfrac{\pi}{2} - \theta\bmod\tfrac{\pi}{2}\bigr).
\]
This is algebraically identical to pole-tracking from \texttt{pole=0} and
runs at C-speed with no Python loop.

\subsection{$\phi$-decomposition: vectorising over $\theta_0$}

With $\omega_0 = 0$ fixed, the analytical trajectory satisfies
$\theta(t;\theta_0) = \theta_0 + \phi(t)$, where $\phi$ is
$\theta_0$-independent.  The grid-search inner loop over $N_{\theta}$
candidates is replaced by a single broadcast:
\begin{verbatim}
phi = _compute_phi(a_I, c_I, t, P, P_1, t_n)   # shape (N,)
theta = theta0_arr[:, None] + phi[None, :]       # shape (N_theta, N)
theta_mod = theta % (pi/2)
d = minimum(theta_mod, pi/2 - theta_mod)
scores = mean((d < k) - sensor, axis=1) ** 2    # shape (N_theta,)
\end{verbatim}
For each $\omega_{ss}$ candidate, $\phi$ is computed once and all
$N_\theta$ scores are evaluated in one pass — no inner Python loop.

\subsection{Hybrid process$+$thread parallelism}

The outer $\omega_{ss}$ grid is split into $n_{\text{proc}}$ chunks
(\texttt{os.cpu\_count()}) and dispatched to a fork-context
\texttt{ProcessPoolExecutor}.  Within each process a
\texttt{ThreadPoolExecutor} runs one thread per $\omega_{ss}$ value.
Under Python~3.14 no-GIL, threads execute NumPy code in true parallel and
share \texttt{sensor\_data}/\texttt{t\_s} by reference, reducing IPC
overhead from $N_\omega$ to $n_{\text{proc}}$ pickle operations.

\subsection{Coarse-to-fine 2D grid}

\textbf{Pass A (coarse):} $100 \times 25 = 2{,}500$ evaluations over
$\omega_{ss,2d} \pm 4\,\text{rad/s} \times [0,\pi/2)$; identifies the
correct phase basin.

\textbf{Pass B (fine):} step $= 10^{-3}$ on both axes.
$\omega_{ss}$: $\pm 0.5\,\text{rad/s}$ around the coarse best ($\approx 1000$ pts).
$\theta_0$: $[\theta_0^* \pm 2k]$ ($\approx 682$ pts).
Total $\approx 682{,}000$ vectorised matrix evaluations.

\subsection{\texttt{P\_SCALE} consistency}

\texttt{generate\_data.py} exposes \texttt{P\_SCALE} (default 1.0) to
simulate at reduced speed.  \texttt{train.py} prompts for the same value
and sets \texttt{P = 1479\,$\times$\,1.05\,$\times$\,P\_SCALE}.
Since $a_I = \omega_{ss}\,c_I/P$, using different \texttt{P} in the two
scripts produces a rescaled but otherwise correct $a_I$
(specifically, the recovered $a_I$ equals the true value times
\texttt{P\_SCALE}).

%--------------------------------------------------------------
\section{Key Findings and Open Issues}
%--------------------------------------------------------------

\begin{enumerate}
  \item \textbf{Data alignment bug (fixed).}  The original \texttt{load\_data()}
        re-zeroed $t$ \emph{before} filtering to $t\geq 2.26\,\text{s}$,
        causing all model evaluations to use the wrong time origin.  The fix
        adds a second \texttt{t = t - t[0]} \emph{after} filtering, aligning
        $t=0$ with PWM onset.  All prior optimisation results (spurious
        $\omega_0 \sim 944\,\text{rad/s}$, $c_I \sim 6$) were artefacts of
        this misalignment.
  \item \textbf{$\omega_0$ lower bound.}  The bound was corrected from
        $(1,1000)$ to $(0,1000)$ to allow the motor-starts-from-rest scenario.
  \item \textbf{BCE / score inversion.}  The soft BCE loss and binary MSE
        have different local minima.  Higher sharpness worsens score because
        confident wrong predictions are heavily penalised, pushing the
        optimizer toward uncertain ($\hat{p}\approx 0.5$) predictions.
  \item \textbf{Optimal sharpness $s^*\approx 4 \approx 1/k$.}  Coupling
        $s=1/k$ removes the sharpness hyperparameter entirely.
  \item \textbf{Best score so far: 0.167} (analytical pipeline +
        coarse-to-fine 2D grid at $P_{\text{scale}}=0.46$,
        $\omega_{ss}=269.6\,\text{rad/s}$; all parameters correct to
        $<1\%$ after \texttt{P\_SCALE} consistency fix).
  \item \textbf{3-sensor quantization floor.}  At $\omega_{ss} > \pi/(6\Delta t)
        \approx 544\,\text{rad/s}$, the 3-sensor estimator is saturated and
        returns the floor value.  The single-sensor estimator is more accurate
        in this regime (transition intervals $> \Delta t$).
  \item \textbf{Phase drift bottleneck.}  Even a 0.017\% error in $\omega_{ss}$
        causes 1\,rad drift over 10\,s, degrading score by $\sim 50\%$.
        Mitigated by the parallel exhaustive grid in Phase 3.
  \item \textbf{$a_I$--$c_I$ scale degeneracy.}
        Any pair $(\lambda a_I, \lambda c_I)$ for $\lambda>0$ gives the same
        $\omega_{ss} = a_IP/c_I$.  The coast-down log-linear fit breaks this
        degeneracy analytically by estimating $c_I$ independently of $a_I$.
  \item \textbf{Open next steps:}
        (a) collect real motor data (single + 3-bit Hall state, \texttt{micros()} timestamps);
        (b) apply the analytical pipeline to real data and compare recovered
            parameters against bench measurements of $\tau = 1/c_I$ and
            no-load $\omega_{ss}$.
\end{enumerate}

\end{document}
